
\documentclass{article}
\usepackage{colortbl}
\usepackage{makecell}
\usepackage{multirow}
\usepackage{supertabular}

\begin{document}

\newcounter{utterance}

\centering \large Interaction Transcript for game `cladder', experiment `full\_v1.5\_default', episode 7019 with refineonly{-}Qwen3.5{-}9B.
\vspace{24pt}

{ \footnotesize  \setcounter{utterance}{1}
\setlength{\tabcolsep}{0pt}
\begin{supertabular}{c@{$\;$}|p{.15\linewidth}@{}p{.15\linewidth}p{.15\linewidth}p{.15\linewidth}p{.15\linewidth}p{.15\linewidth}}
    \# & \multicolumn{2}{c}{Player} && \multicolumn{2}{c}{Game Master} \\
    \hline

    \theutterance \stepcounter{utterance}  
    & & & \multicolumn{4}{p{0.6\linewidth}}{
        \cellcolor[rgb]{0.9,0.9,0.9}{
            \makecell[{{p{\linewidth}}}]{
                \texttt{\tiny{[P1$\langle$GM]}}
                \texttt{You are an expert in causal inference. The following question is not a typical commonsense query, but rather a meticulously designed question created by a professor specializing in causal inference, intended to assess the students' mastery of the course content.} \\
\texttt{Imagine a self{-}contained, hypothetical world with only the following conditions, and without any unmentioned factors or causal relationships: Hwax has a direct effect on gyzp and jyka. Jyka has a direct effect on lirg. Gyzp has a direct effect on lirg. Method 1: We look at how jyka correlates with lirg case by case according to gyzp. Method 2: We look directly at how jyka correlates with lirg in general. To understand how jyka affects lirg, is it more correct to use the Method 1 than Method 2?} \\
\texttt{Start your answer with 'Yes' or 'No', followed by additional reasoning or evidence to support your explanation.} \\
            }
        }
    }
     \\ \\

    \theutterance \stepcounter{utterance}  
    & \multicolumn{4}{p{0.6\linewidth}}{
        \cellcolor[rgb]{0.561,0.737,0.561}{
            \makecell[{{p{\linewidth}}}]{
                \texttt{\tiny{[P1$\rangle$GM]}}
                \texttt{No} \\
\\ 
\texttt{**Reasoning:**} \\
\\ 
\texttt{In this hypothetical world, the causal structure forms a scenario involving a **confounder** (\$Hwax\$) and a **collider** (\$lirg\$) relative to the path of interest, but more critically, it presents a classic case of **confounding bias** in the marginal association.} \\
\\ 
\texttt{Let's map the graph:} \\
\texttt{1.  \$Hwax \textbackslash{}rightarrow gyzp\$} \\
\texttt{2.  \$Hwax \textbackslash{}rightarrow jyka\$} \\
\texttt{3.  \$gyzp \textbackslash{}rightarrow lirg\$} \\
\texttt{4.  \$jyka \textbackslash{}rightarrow lirg\$} \\
\\ 
\texttt{Here, \$Hwax\$ is a common cause (confounder) of both the exposure (\$jyka\$) and a variable (\$gyzp\$) that also affects the outcome (\$lirg\$). This creates a backdoor path: \$jyka \textbackslash{}leftarrow Hwax \textbackslash{}rightarrow gyzp \textbackslash{}rightarrow lirg\$.} \\
\\ 
\texttt{If we look at **Method 2** (the general/marginal correlation), we are observing the total association between \$jyka\$ and \$lirg\$. This association is biased because it includes:} \\
\texttt{1.  The direct causal effect of \$jyka\$ on \$lirg\$.} \\
\texttt{2.  The spurious association induced by the confounder \$Hwax\$ flowing through \$gyzp\$.} \\
\texttt{Therefore, Method 2 does not isolate the causal effect of \$jyka\$.} \\
\\ 
\texttt{If we look at **Method 1** (the case{-}by{-}case correlation according to \$gyzp\$), we are stratifying by \$gyzp\$. By holding \$gyzp\$ constant, we block the backdoor path \$jyka \textbackslash{}leftarrow Hwax \textbackslash{}rightarrow gyzp \textbackslash{}rightarrow lirg\$. Within each stratum of \$gyzp\$, the variation in \$lirg\$ is no longer confounded by \$Hwax\$ via \$gyzp\$. Consequently, the correlation observed within these strata isolates the direct causal effect of \$jyka\$ on \$lirg\$.} \\
\\ 
\texttt{Thus, to understand the true causal effect of \$jyka\$ on \$lirg\$ and remove the bias introduced by the common cause \$Hwax\$, it is indeed more correct to use Method 1 (stratification) than Method 2 (marginal correlation).} \\
            }
        }
    }
    & & \\ \\

    \theutterance \stepcounter{utterance}  
    & & & \multicolumn{2}{p{0.3\linewidth}}{
        \cellcolor[rgb]{0.9,0.9,0.9}{
            \makecell[{{p{\linewidth}}}]{
                \texttt{\tiny{[GM$|$GM]}}
                \texttt{No} \\
            }
        }
    }
    & & \\ \\

    \theutterance \stepcounter{utterance}  
    & & & \multicolumn{2}{p{0.3\linewidth}}{
        \cellcolor[rgb]{0.9,0.9,0.9}{
            \makecell[{{p{\linewidth}}}]{
                \texttt{\tiny{[GM$|$GM]}}
                \texttt{yes} \\
            }
        }
    }
    & & \\ \\

    \theutterance \stepcounter{utterance}  
    & & & \multicolumn{2}{p{0.3\linewidth}}{
        \cellcolor[rgb]{0.9,0.9,0.9}{
            \makecell[{{p{\linewidth}}}]{
                \texttt{\tiny{[GM$|$GM]}}
                \texttt{game\_result = LOSE} \\
            }
        }
    }
    & & \\ \\

\end{supertabular}
}

\end{document}
